% !PRINTED 175
\weberpage{175}

\begin{center}
§. 54.

\gesperrt{Die Tschirnhausen-Transformation der Gleichung fünften Grades.}
\end{center}

Will man auf die Gleichung fünften Grades
\weq{1}{x^5 + a_1 x^4 + a_2 x^3 + a_3 x^2 + a_4 x + a_5 = 0}
die Tschirnhausen-Transformation anwenden, so hat man
\weq{2}{y = \alpha_0 + \alpha_1 x + \alpha_2 x^2 + \alpha_3 x^3 + \alpha_4 x^4}
zu setzen, und erhält die Transformirte
\weq{3}{y^5 + b_1 y^4 + b_2 y^3 + b_3 y^2 + b_4 y + b_5 = 0.}

Der nächstliegende Gedanke wäre nun, die Verhältnisse der
fünf Unbekannten $\alpha$ so zu bestimmen, dass
\weq{4}{b_1 = 0, \quad b_2 = 0, \quad b_3 = 0, \quad b_4 = 0}
würde, wodurch (3) auf eine reine Gleichung zurück käme, und
dies war wohl auch das Ziel, das \gesperrt{Tschirnhausen} im Auge
hatte. Aber von den vier Gleichungen (4) ist die erste linear,
die zweite quadratisch, die dritte vom dritten und die vierte vom
vierten Grade. Die Endgleichung, die man daraus ableiten kann,
wird daher (nach \S.\,50) vom Grade $2.3.4 = 24$, und daher
kann daraus für die Lösung der Gleichung fünften Grades kein
Nutzen gezogen werden. Man sucht daher zunächst nur den
beiden Gleichungen zu genügen
\weq{5}{b_1 = 0, \quad b_2 = 0,}
wodurch die Gleichung (3) in eine Form übergeht, die nach
\gesperrt{F. Klein} (Vorlesungen über das Ikosaëder) eine \gesperrt{Hauptgleichung} fünften Grades genannt wird.

Zur Befriedigung der beiden Gleichungen (5) haben wir die
Verfügung über die vier Verhältnisse $\alpha_0 : \alpha_1 : \alpha_2 : \alpha_3 : \alpha_4$, und wir
können diese Grössen noch auf mannigfaltige Art zur Vereinfachung der Gleichung fünften Grades benutzen.

Wir kommen in einem späteren Abschnitt ausführlicher auf
diesen Gegenstand zurück und beschränken uns hier darauf, die
Ziele im Allgemeinen zu bezeichnen. Der Anschaulichkeit wegen
bedienen wir uns einer geometrischen Einkleidung, die übrigens
zum Verständniss der algebraischen Theorie, wie sie später
gegeben werden soll, durchaus nicht wesentlich ist.
